% Active Figure 2 caption; mirrored in proposal.tex.
\caption{\textbf{The spectral, spatial, and temporal ingredients of the experiment, each already shown to work in a single AGN.}
\textbf{Top---Spectral: measuring the dust spectrum.} The NIRSpec+MIRI/MRS spectrum of
NGC~7469 (digitized from [15, Fig.~8]) separates cleanly into hot and warm dust
components. The shaded 8--13~$\mu$m band provides the dust luminosity once PAH
and line emission are subtracted, and the 9.7 and 18~$\mu$m silicate features
add a spectral-shape diagnostic that imaging cannot supply. Bars mark MRS and,
for comparison, SPHEREx coverage at $z=0.0163$.
\textbf{Lower left---Spatial: isolating the nucleus.} Spectra and 11.6~$\mu$m maps of
NGC~6552 [28, Fig.~4] decompose the emission into nuclear (PSF) and
circumnuclear components, the step that decides whether a nuclear warm excess
survives host subtraction.
\textbf{Lower right---Temporal: reading the dust's memory.} In NGC~4151 the 10.5~$\mu$m
emission (red points; light-blue curves are fitted responses) follows the
optical light curve (blue) with an $\sim$8~yr delay and additional smoothing
[23, Fig.~15], preserving a record of earlier nuclear activity.
Wavelengths are rest-frame in the top panel and observed-frame at lower left;
the lower-right panel omits longer-wavelength data and ancillary fit labels.
Our programme is the first to combine all three steps across a sample---24 nuclei
with recorded rises, declines, flares, and recoveries---and so to test where
delayed heating alone explains the mid-infrared response and where the dust
distribution itself must evolve.}
